\section{Method: $p$-fair BPE}
\label{sec:method}

Like classical BPE, $p$-fair BPE builds the vocabulary one merge at a time and chooses each
merge greedily. We first derive the merge rule implied by Eq.~\eqref{eq:obj}
(Section~\ref{sec:rule}), then show that it contains classical and Parity-aware BPE as
limits (Section~\ref{sec:limits}), and finally give an exact and fast training algorithm
(Section~\ref{sec:lazy}).

\subsection{The greedy merge rule}
\label{sec:rule}

Merging a pair $(a,b)$ replaces each of its occurrences by one symbol, so the token count
of a text falls by the number of occurrences of the pair. The pair counts $c_\ell$ come from
the training text, whereas the costs $t_\ell$ are measured on the development set. To
connect the two, we assume that a pair occurs in the development text of language $\ell$ at
the same rate as in its training text. Let $N_\ell$ denote the size of $\mathcal{D}_\ell$
expressed in development sentences, that is, its length divided by the average length of a
development sentence in language $\ell$. Under this assumption, merging $(a,b)$ lowers
$T_\ell$ by $S\,c_\ell(a,b)/N_\ell$ and therefore lowers $t_\ell$ by $c_\ell(a,b)/N_\ell$.

\begin{proposition}
\label{prop:rule}
Let $\pi_\ell = N_\ell/\sum_k N_k$. To first order, the merge that decreases $F_p$ the most
is
\begin{equation}
  (a,b)^\star \;=\; \argmax_{(a,b)}\; s_p(a,b), \qquad
  s_p(a,b) = \sum_{\ell\in\mathcal{L}} t_\ell^{\,p-1}\, c_\ell(a,b).
  \label{eq:rule}
\end{equation}
\end{proposition}

\begin{proof}
Since $x\mapsto x^{1/p}$ is increasing, minimising $F_p$ is equivalent to minimising
$\Phi_p=\sum_\ell \pi_\ell\, t_\ell^{\,p}$. A merge of $(a,b)$ changes each $t_\ell$ by
$\Delta t_\ell = -c_\ell(a,b)/N_\ell$, so the first-order change in $\Phi_p$ is
\begin{equation*}
  \Delta\Phi_p \approx \sum_\ell p\,\pi_\ell\, t_\ell^{\,p-1}\Delta t_\ell
  = -\frac{p}{\sum_k N_k}\sum_\ell t_\ell^{\,p-1} c_\ell(a,b),
\end{equation*}
where the factors $N_\ell$ cancel with $\pi_\ell$. The decrease is largest for the pair
that maximises $s_p$.
\end{proof}

Thus, the rule is a pair count in which each language is weighted by its current cost raised
to the power $p-1$. The weights change after every merge, since each merge lowers the costs
of the languages in which the pair occurs. We update every $t_\ell$ exactly after each merge
by applying the merge to the development set, so only the per-merge choice relies on the
first-order approximation.

To illustrate the effect of the weights, under classical BPE in our experiments a Hindi
sentence costs 35.0 tokens and its English translation 29.0. A Hindi pair then counts
$(35.0/29.0)^{p-1}$ times as much as an English pair of the same frequency: 1.2 times at
$p=2$, 1.8 times at $p=4$ and 3.7 times at $p=8$.

\subsection{Classical and Parity-aware BPE as limits}
\label{sec:limits}

The rule in Eq.~\eqref{eq:rule} has two limits.
\begin{itemize}
  \item \textbf{$p=1$.} Every weight equals one, $s_1(a,b)=\sum_\ell c_\ell(a,b)$ is the
  count of the pair in the pooled training text, and the rule is exactly classical BPE.
  \item \textbf{$p\to\infty$.} Dividing $s_p$ by $(\max_\ell t_\ell)^{p-1}$ does not
  change the $\argmax$. The weight of the worst-off language is then one and every other
  weight tends to zero, so the rule selects the most frequent pair of the worst-off
  language. This is the base rule of Parity-aware BPE \citep{foroutan2026} with cost
  measured per parallel sentence, and corresponds to a subgradient step on
  $\max_\ell t_\ell$.
\end{itemize}
For $1<p<\infty$, every language keeps a positive weight. In contrast to the parity rule,
which considers one language per merge, and to the Worst-$N$ rule \citep{kanjirangat2026},
which weights the $N$ worst-off languages equally, $p$-fair BPE weights all languages
continuously by their cost. In contrast to the hybrid schedule \citep{foroutan2026}, which
switches rules at a fixed merge index, the balance between compression and parity is set by
one parameter that is constant throughout training.

\subsection{Exact training by lazy evaluation}
\label{sec:lazy}

Efficient implementations of classical BPE keep the pair counts in a priority queue and
update only the pairs that a merge touches. In Eq.~\eqref{eq:rule} the weights change after
every merge, so every score changes, and rescoring all pairs at each of the
$\sim$31{,}000 steps costs $O(|\mathcal{P}|\,L)$ per step, where $\mathcal{P}$ is the set of
pairs. To avoid this, we use the following property.

\begin{proposition}
\label{prop:lazy}
For $p\ge1$, the score $s_p(a,b)$ of a pair that already exists never increases during
training.
\end{proposition}

\begin{proof}
Adding a merge to the vocabulary can only join symbols, so no text gains tokens and every
$t_\ell$ is non-increasing. Since $p-1\ge0$, every weight $t_\ell^{\,p-1}$ is
non-increasing as well. A merge of $(a,b)$ into a new symbol $z$ removes occurrences of
pairs that overlap $(a,b)$ and creates new adjacencies only with $z$, so the count
$c_\ell$ of any pair that existed before the merge can only fall. The score $s_p$ is a sum of
products of non-negative, non-increasing terms, and is therefore non-increasing.
\end{proof}

Consequently, a stored score is always an upper bound on the current score. We keep the
pairs in a max-heap keyed by their stored scores. At each step, we pop the top entry and
recompute its score with the current weights. If the score is unchanged, the entry is the
exact maximiser of Eq.~\eqref{eq:rule}, since every other stored score bounds its current
score from above; otherwise, we push it back with the new score and repeat
(Algorithm~\ref{alg:pfair}). Pairs that contain a new symbol enter the heap with their
exact scores. This is the lazy greedy technique of \citet{minoux1978}, and it returns the
same merge as full rescoring at every step. For $p=\infty$, we keep one heap per language
keyed by $c_\ell$ and pop from the heap of the worst-off language.

\begin{algorithm}[t]
\caption{$p$-fair BPE training with lazy evaluation.}
\label{alg:pfair}
\small
\begin{algorithmic}[1]
\Require training texts $\{\mathcal{D}_\ell\}$, development set, exponent $p$, budget $K$
\State $V \gets$ base alphabet; compute $c_\ell(a,b)$ and $t_\ell$ for all $\ell$
\State $H \gets$ max-heap of all pairs keyed by $s_p(a,b)$
\While{$|V| < K$}
  \State pop $(a,b)$ with stored score $\hat s$ from $H$
  \State $s \gets \sum_\ell t_\ell^{\,p-1} c_\ell(a,b)$ \Comment{current score}
  \If{$s < \hat s$}
    \State push $(a,b)$ with score $s$; \textbf{continue}
  \EndIf
  \State add merge $(a,b)\to z$ to $V$; update $c_\ell$ for touched pairs
  \State push pairs containing $z$ with their exact scores
  \State apply the merge to the development set; update every $t_\ell$
\EndWhile
\State \Return $V$
\end{algorithmic}
\end{algorithm}